% Options for packages loaded elsewhere
\PassOptionsToPackage{unicode}{hyperref}
\PassOptionsToPackage{hyphens}{url}
%
\documentclass[
  11pt,
]{article}
\usepackage{amsmath,amssymb}
\usepackage{iftex}
\ifPDFTeX
  \usepackage[T1]{fontenc}
  \usepackage[utf8]{inputenc}
  \usepackage{textcomp} % provide euro and other symbols
\else % if luatex or xetex
  \usepackage{unicode-math} % this also loads fontspec
  \defaultfontfeatures{Scale=MatchLowercase}
  \defaultfontfeatures[\rmfamily]{Ligatures=TeX,Scale=1}
\fi
\usepackage{lmodern}
\ifPDFTeX\else
  % xetex/luatex font selection
\fi
% Use upquote if available, for straight quotes in verbatim environments
\IfFileExists{upquote.sty}{\usepackage{upquote}}{}
\IfFileExists{microtype.sty}{% use microtype if available
  \usepackage[]{microtype}
  \UseMicrotypeSet[protrusion]{basicmath} % disable protrusion for tt fonts
}{}
\makeatletter
\@ifundefined{KOMAClassName}{% if non-KOMA class
  \IfFileExists{parskip.sty}{%
    \usepackage{parskip}
  }{% else
    \setlength{\parindent}{0pt}
    \setlength{\parskip}{6pt plus 2pt minus 1pt}}
}{% if KOMA class
  \KOMAoptions{parskip=half}}
\makeatother
\usepackage{xcolor}
\usepackage[margin=1in]{geometry}
\usepackage{listings}
\newcommand{\passthrough}[1]{#1}
\lstset{defaultdialect=[5.3]Lua}
\lstset{defaultdialect=[x86masm]Assembler}
\setlength{\emergencystretch}{3em} % prevent overfull lines
\providecommand{\tightlist}{%
  \setlength{\itemsep}{0pt}\setlength{\parskip}{0pt}}
\setcounter{secnumdepth}{-\maxdimen} % remove section numbering
\lstset{basicstyle=\ttfamily\scriptsize,breaklines=true,columns=fullflexible,keepspaces=true}
\emergencystretch=2em
\ifLuaTeX
  \usepackage{selnolig}  % disable illegal ligatures
\fi
\usepackage{bookmark}
\IfFileExists{xurl.sty}{\usepackage{xurl}}{} % add URL line breaks if available
\urlstyle{same}
\hypersetup{
  pdftitle={Completed algebraic zero-Hessian ascent on pure analytic curves},
  pdfauthor={Xean research manuscript},
  hidelinks,
  pdfcreator={LaTeX via pandoc}}

\title{Completed algebraic zero-Hessian ascent on pure analytic curves}
\author{Xean research manuscript}
\date{3 October 2026}

\begin{document}
\maketitle

\subsection{1. Exact result and logical
scope}\label{exact-result-and-logical-scope}

Put \(D=d'd''\) and \(t_f=-\log|f|\). All smooth forms are
\textbf{ordinary CLD forms}, represented on ordinary open neighbourhoods
by analytic unit moments. All measures use the canonical calibration of
the actual analytic algebra, including its Artinian lengths. No space is
replaced by its reduction for integration.

\textbf{Theorem (completed algebraic curve ascent).} Let \(K\) be
complete and nontrivially valued, let \(M/K\) be an algebraic extension
with an extending absolute value, and let \(L=\widehat M\). Let \(X\) be
a good, separated, strict, boundaryless, pure one-dimensional
\(K\)-analytic space. Singular and nonreduced spaces are allowed. For
every continuous finite real function \(h\) and the scalar projection
\(\pi:X_L\to X\), \[
             D[h]=0\quad\Longleftrightarrow\quad
             D[h\circ\pi]=0.                         \tag{1}
\] The statement holds on every ordinary open subspace of \(X\). In
particular it holds for the completion of an algebraic closure, for
completed separable algebraic extensions, and in positive characteristic
for arbitrary mixtures of inseparability and separability. \(L/K\)
itself need not be algebraic. No upstairs current equality is imposed.

\textbf{Exact compact scalar test descent.} For a good separated
\(K\)-analytic \(X\), with the same fields, every
\(\beta\in A_c^{0,0}(X_L)\) is the exact smooth pullback of an
\(\alpha\in A_c^{0,0}(X_E)\) for some finite \(K\subset E\subset M\).
Here \(q_E:X_L\to X_E\) is the actual scalar projection and
\(\beta=q_E^*\alpha\) as ordinary smooth scalar sections, so also
\(D\beta=q_E^*D\alpha\). This subsidiary result needs no pure dimension
or boundarylessness. Its proof, including domains, ordinary open loci,
units and compact supports, is in Section 12. It is local descent to a
finite stage; it does not assert descent of every analytic function or
an arbitrary affinoid presentation.

The only approximation needed at the completed stage is of finitely many
analytic functions in Banach norm. Strict ultrametric inequalities turn
it into \textbf{exact} domain and unit-norm identities. No
differentiated limit or infinite group average is taken at that stage.
For finite separable ascent, the separate approximation assertion proved
in Sections 3--10 is: for invariant \(\beta\in A_c^{0,0}(X_E)\) there
are \(\alpha_j\in A_c^{0,0}(X)\) and a fixed compact \(C\subset X\) with
\[
 \operatorname{supp}\alpha_j\subset C,\qquad
 \big\|(D(\beta-\pi_E^*\alpha_j))_{\mathrm{Rad},E}\big\|_{\mathrm{TV}}
                 \longrightarrow0,                  \tag{2}
\] and every error supported in \(\pi_E^{-1}C\). Both approximation
limits, including collision-contained calibrated segments, are
reconstructed. No earlier draft theorem, review, hash or native note is
a premise.

On a pure curve the current equation is precisely \[
 B_h(\psi):=\int_X h\,(D\psi)_{\mathrm{Rad},K}=0
                  \quad(\psi\in A_c^{0,0}(X)).         \tag{3}
\] Only the differentiated smooth top form is represented by a measure;
\(D[h]\) is not presumed to be a Radon measure. Current-duality signs do
not affect vanishing.

This is a scalar transport theorem, a partial result for original P06.2.
It gives no recognition of \(h\), no two-sign CLD-psh regularization, no
universal curve-probe transport, and no higher-dimensional
separable-current ascent. P05.2 is not addressed. The original status is
\passthrough{\lstinline!partial!} and
\passthrough{\lstinline!candidate=false!}.

Sections 2--10 reconstruct the finite separable dependency. Section 11
reconstructs finite Frobenius descent and the actual finite mixed-field
decomposition. Section 12 proves the completion and support gate.
Section 13 concludes the current theorem. Section 15 includes the used
the source audit and the corrected primary-source supplement accompanying this paper.

\subsection{2. Exact analytic dependencies and source
checks}\label{exact-analytic-dependencies-and-source-checks}

CLD refers to Chambert-Loir--Ducros, \emph{Formes différentielles
réelles et courants sur les espaces de Berkovich}, arXiv:1204.6277v2, 25
July 2025. GR refers to Gubler--Rabinoff, \emph{Classically psh and
pluriharmonic functions on Berkovich spaces}, arXiv:2506.13548v2, 17
September 2025. Numbering is that of these cached primary snapshots,
whose statements and proof bodies, not just the former drafts'
descriptions, were inspected.

\begin{enumerate}
\def\labelenumi{\arabic{enumi}.}
\item
  \textbf{Affinoid foundations.} CLD 11.1.1--2: an affinoid spectrum is
  compact Hausdorff; its algebra is noetherian with closed ideals; the
  map to its prime spectrum is surjective. CLD 11.1.5: every point of an
  affinoid has rational affinoid neighbourhoods inside any prescribed
  ordinary open, with the point satisfying all associated strict
  inequalities. Rational systems have positive real radii and
  denominator/numerators with no common zero (Section 12 uses this
  stronger neighbourhood clause of 11.1.5). CLD 11.1.7--8: a good space
  has an affinoid neighbourhood basis; ordinary opens are analytic
  domains; analytic sections form a sheaf also for ordinary opens, and
  \(\mathcal O(\mathcal M(A))=A\). CLD 11.1.9: affinoid fibre products
  and scalar extensions use completed tensor products, and scalar
  projection is surjective. CLD 11.4.4: a finite algebra with its
  finite-module Banach topology is affinoid. CLD 11.7.7: over a
  nontrivially valued field, every point of a good strict space has a
  strict affinoid neighbourhood. Applying this to ordinary opens gives
  such a neighbourhood inside the open. Particular tensor, orbit,
  compactness, shrinking and unit claims are proved in Section 3.
\item
  \textbf{Dimension and boundary.} CLD 11.2.2: separated analytic spaces
  are topologically separated. CLD 11.2.3: strict affinoid dimension is
  Krull dimension, dimension is preserved by complete scalar extension,
  and domains preserve local analytic dimension. If the minimal-prime
  components of \(\mathcal M(A)\) are \(X_i\), then
  \(\dim_xX=\sup_{x\in X_i}\dim X_i\). Pure dimension means constant
  local dimension. GR 2.11(5), equivalently (5'), states for a morphism
  \(Y\to Z\) and any analytic extension \(L/K\) that
  \(\partial(Y_L/Z_L)\subseteq\pi^{-1}(\partial(Y/Z))\). Hence a
  boundaryless space remains boundaryless. GR 2.11(3),(4) gives
  boundarylessness for an ordinary open in a boundaryless space: its
  relative interior is its ordinary topological interior and interiors
  compose. Section 3 proves purity component by component; equality of
  total dimensions alone is not the proof.
\item
  \textbf{Ordinary forms.} CLD Corollary 15.1.2 states that a function
  is smooth near \(x\) if it has a representation
  \(\varphi(t_{f_1},\ldots,t_{f_m})\) on a neighbourhood of \(x\), where
  the \(f_i\) are analytic units and \(\varphi\) is smooth on a
  Euclidean neighbourhood of the coordinate value. CLD Proposition
  15.1.1(b) gives ordinary open moment descriptions for forms in every
  bidegree. CLD 6.2.1--2 defines ordinary forms and says they form an
  ordinary sheaf stable under wedge, \(J,d',d''\). Moment pullback
  commutes with the exterior calculus (CLD 3.3.12). In a moment a
  logarithm is an affine coordinate, so \(Dt_f=0\) as a \textbf{form}.
  CLD 3.4.4: for a G-tropical chart \(f\) and a coefficient form \(a\),
  its pullback is zero in the G-form sheaf exactly when \(a\) restricts
  to zero on \(f(V)\) for every compact domain \(V\). This identifies
  the ordinary form sheaves when the analytic units, domains and moment
  images are the same; it is not descent through scalar projection.
\item
  \textbf{Compact cutoffs only.} CLD 15.2.2: a good topologically
  separated space is locally holomorphically separated. CLD Corollary
  15.2.5: on a locally holomorphically separated analytic space, for any
  compact \(F\) and ordinary open neighbourhood \(U\) of \(F\), there is
  an ordinary smooth function equal to one near \(F\) with compact
  support contained in \(U\). Its proof uses local differential
  separation (15.2.4) and the smooth bump result 6.2.4. Section 4
  constructs the only partition statement needed, for compact \(F\) and
  finite covers. The overbroad fixed-index countable-at-infinity reading
  of CLD 15.2.9 is not used.
\item
  \textbf{Canonical top measures.} CLD Theorem 15.7.1: a pure
  \(n\)-dimensional topologically separated analytic space with its
  canonical calibration and analytic boundary is a tropical space of
  level \(n\). CLD 15.7.2: the induced structure on a domain agrees with
  that domain's analytic structure. CLD 5.4.4: every compact tropical
  chart is pleasant, and its tropical polyhedron is the top-dimensional
  part of its compact moment image. CLD 14.3.2--5, 14.5.1--2: canonical
  calibrations are standard and positive, and on a top moment cell \(Q\)
  the calibration is \(d_Q\tau_{Q,*}|e_1\wedge\cdots\wedge e_n|\), where
  \(\tau_Q\) is the affine inverse of an injective torus projection and
  \(d_Q\) is the sum of the actual finite-flat degrees of the skeletal
  cells above \(Q\). Those degrees retain fibre-algebra dimensions,
  including lengths. CLD 5.8.3--5 constructs a top form's measure by
  contraction on its supporting skeleton. CLD Proposition 5.10.1(c,d):
  for a pleasant chart \(f\) and a Borel \((n,n)\) coefficient form
  \(a\), the induced measure is locally finite exactly when \(a\) is
  locally integrable on the calibrated tropical polyhedron, and then
  \(f_*((f^*a)_{\mathrm{Rad}})=\langle a,\mu_{\Pi_f}\rangle_{\mathrm{Rad}}\);
  the corresponding integrals agree. In particular a smooth top form is
  finite on a compact moment domain. Lower-dimensional cells have no top
  measure. Compact moment images are finite polyhedra by the
  polyhedrality assertion of CLD Theorem 12.2.4. The dimension bound is
  specifically CLD 12.2.10(1)--(2):
  \(\dim_{\mathrm{trop},x}f\le d_K(x)\le\dim_xX\) and, for compact
  \(X\), \(\dim f_{\mathrm{trop}}(X)=\sup_x\dim_{\mathrm{trop},x}f\).
  Thus compact curve images have dimension at most one. We do not use a
  claim that a moment image always has dimension equal to its source;
  constant moments rule out that reading. CLD 7.3.5 defines integration
  currents for locally integrable multipliers; a continuous finite
  multiplier is locally integrable. Section 8 proves the carrier
  total-variation estimate needed here, before any possibly cancelling
  signed pushforward.
\item
  \textbf{Scalar projection has no field-degree factor.} CLD Proposition
  14.3.6: on a level-\(n\) analytic space, if \(\Sigma\) has canonical
  calibration \(\mu\) and \(\Sigma_L\) its scalar extension has
  canonical calibration \(\mu_L\), then \(\pi_*\mu_L=\mu\). CLD
  Proposition 15.3.1: for a tropically Borel \((n,n)\) G-form
  \(\omega\), \(\pi_*\omega_{L,\mathrm{Bor}}=\omega_{\mathrm{Bor}}\);
  for a locally integrable such form, \[
             \pi_*\omega_{L,\mathrm{Rad}}
                       =\omega_{\mathrm{Rad}}.          \tag{4}
  \] Here is the primary proof mechanism, retaining its lengths. On a
  convenient skeletal cell choose a torus projection injective on the
  cell. Its finite-flat degree downstairs equals the sum of its degrees
  at all upstairs skeletal lifts. This is CLD Lemma 13.3.3(a): degrees
  at \(y\) equal the sum of degrees over \(y\) after complete base
  change and choice of a target lift. The lemma reduces to a fibre; a
  zero-dimensional local fibre algebra is Artinian, and its tensor
  decomposes into its local Artinian factors. The sum of their
  dimensions is the original dimension after scalar extension.
  Positive-dimensional fibres give infinite degrees on both sides. Thus
  the cell volume vectors sum to the downstairs vector, proving 14.3.6.
  Contracting the form and taking Radon measures proves 15.3.1. We apply
  (4) only to smooth top forms, then pair with \(h\); this never assumes
  pullback of a continuous current.
\item
  \textbf{Same-carrier field restriction has a degree factor.} CLD
  Remark 14.4.3: for a finite extension \(F/K\), an \(F\)-analytic space
  \(Z\) of level \(n\), and its restricted \(K\) structure \(Z_0\),
  dimensions and skeletal subsets agree and \(\mu_{Z_0}=[F:K]\mu_Z\).
  The proof uses the finite-flat degree \([F:K]\) of the projection of
  the split torus on this restricted carrier. CLD 15.4.1 and Proposition
  15.4.3: for a morphism \(f:Y\to Z\) of level-\(n\) spaces finite off a
  Zariski closed set of dimension \(<n\), a locally integrable \((p,n)\)
  G-form \(\omega\), \(p\in\{n-1,n\}\), has locally integrable pullback
  and
  \(f_*((f^*\omega)_{\mathrm{Rad}})=\mathrm{dc}(f)\omega_{\mathrm{Rad}}\).
  Its proof first proves the calibrated volume-vector identity on the
  inverse skeletal set (14.4.2), then contracts and integrates. We use
  only an actual isomorphism on \(Z_0\), whose calibrated degree is one:
  both local Artinian rings are identical. Cancelling the same-carrier
  factor \([F:K]\) proves invariance over \(F\).
\end{enumerate}

All these source statements apply to the actual, potentially singular
and nonreduced structure. They require no normalization and no
quasi-smoothness. No weakly smooth GJR descent, positive-current
recognition, or catalogue theorem is used as an ordinary form theorem.

\subsubsection{2.1. Direct proofs of the neighbourhood and
compact-cutoff facts used
here}\label{direct-proofs-of-the-neighbourhood-and-compact-cutoff-facts-used-here}

The finite proof uses the primary foundations above. For the two
topological localization facts decisive at completion, the following
direct constructions supply their proof bodies as well.

On \(\mathcal M(A)\) the ordinary topology is initial for all
evaluations \(|a|\) (CLD 11.1.1). A prescribed neighbourhood of a point
therefore contains a finite intersection of conditions
\(|a_j|>\ell_j>0\) and \(|b_k|<u_k\), where thresholds can be chosen
strictly between the point's value and the excluded endpoint. Missing
lower or upper bounds impose no condition. Put \(g=\prod_j a_j\), with
\(g=1\) when there are no lower bounds. Choose \(C>0\) so that
\(C^{-1}<\prod_j\ell_j\). The rational conjunction \[
 |1|\le C|g|,\qquad
 |\prod_{j'\ne j}a_{j'}|\le\ell_j^{-1}|g|\ (\forall j),\qquad
 |b_kg|\le u_k|g|\ (\forall k)                         \tag{R}
\] has no common zero among its denominator and numerators, because one
numerator is \(1\). It is a rational analytic domain by the universal
rational-domain construction of CLD 11.1.5. Its validity locus lies in
the originally prescribed neighbourhood if the thresholds were chosen
inside that neighbourhood's intervals: its first inequality excludes
\(g=0\), so division of norms gives all chosen lower and upper bounds.
The original point satisfies every strict inequality, including the
first. Non-strict thresholds in (R) were chosen inside the original open
intervals, so validity of the weak system suffices for containment. This
proves the strong rational-neighbourhood clause needed in Section 12
using the initial topology and the rational-domain construction. Real
positive thresholds suffice even when the value group is discrete.

Here is also a direct compact-cutoff construction in every good
Hausdorff analytic space occurring in the proof. For a point \(x\) in an
ordinary open \(U\), take an affinoid neighbourhood \(V_0\) and within
its ordinary interior a rational neighbourhood \(R\) of \(x\) contained
in \(U\), with strict inequalities \(|f_i(x)|<r_i|g(x)|\). Its
denominator satisfies \(g(x)\ne0\). Choose \(b_i\) with
\(|f_i(x)|/|g(x)|<b_i<r_i\) and choose \(0<a<|g(x)|\). On the ordinary
nonzero locus of \(g\) in \(V_0\), a smooth function of \(\log|g|\)
equals one near \(x\) and vanishes where \(|g|\le a\). For each \(i\),
choose a smooth real function of \(\log|f_i/g|\) that is one when
\(|f_i/g|\le b_i\) and zero when \(|f_i/g|\ge r_i'\) for some
\(b_i<r_i'<r_i\). Extend that factor as the constant one at zeros of
\(f_i\): near such a zero its ratio is below \(b_i\), so the extension
is locally constant and ordinary smooth. On its nonzero locus the ratio
is an analytic unit, giving the required ordinary moment description.
Multiply the factors with the denominator cutoff. At points where
\(g=0\), extend by zero; the denominator cutoff already vanishes on a
neighbourhood of each such point.

The support of this product lies in the compact set in \(V_0\) defined
by \(|g|\ge a\) and \(|f_i|\le r_i'|g|\) for all \(i\). That compact set
lies in the strict locus of \(R\), hence in
\(U\cap\operatorname{int}_X(V_0)\). Extend the product by zero outside
\(V_0\): every point outside the compact support has a zero
neighbourhood, so the extension is ordinary smooth. It equals one near
\(x\). Taking squares gives nonnegative cutoffs. For compact
\(F\subset U\), finitely many positive loci of these point cutoffs cover
\(F\). Their sum \(s\) is positive near \(F\). Compose with a real
smooth function that vanishes near zero and equals one on a
neighbourhood of the compact positive set \(s(F)\). This gives a compact
cutoff equal to one near \(F\) and supported in \(U\). This reconstructs
the compact CLD 15.2.5 conclusion in our actual good Hausdorff scope.
Section 4 then proves its finite-cover partition specialization. No
countable-at-infinity fixed-index compact-support assertion is used.

In particular all cutoffs are obtained from ordinary analytic-unit
moments, flat constant extensions near analytic zeros and compact
extension by zero. No weakly smooth tropicalization is substituted.

\subsection{3. Finite fields, fibres, analytic structures, and
purity}\label{finite-fields-fibres-analytic-structures-and-purity}

We first treat a finite Galois extension \(E/K\), with group \(G\) and
\(d=[E:K]=|G|\). The elementary finite field facts used below can be
justified as follows. A finite separable extension has as many distinct
embeddings in an algebraic closure as its degree: extend embeddings
successively by the distinct roots of each separable minimal polynomial.
The field \(K\) is infinite. For two separable generators \(u,v\), all
but finitely many \(c\in K\) make the values of \(u+cv\) distinct for
every pair of distinct embeddings; if their \(v\) values agree their
\(u\) values already differ. Its minimal polynomial then has at least as
many roots as the extension degree, so \(u+cv\) generates the extension.
Induction gives a primitive element. The splitting field of its
separable minimal polynomial is finite separable. Every embedding of
this splitting field permutes the roots and is onto, so it is Galois and
its automorphism group has its degree as order. Its fixed field is
\(K\): the embeddings of any intermediate separable simple field extend
successively to the splitting field, so a fixed element has a minimal
polynomial of degree one. This also provides the normal closure used in
Section 10.

Finite-dimensional norms over a complete non-Archimedean field are
equivalent to the coefficient maximum norm in a fixed basis. Here is a
proof. Induct on dimension. The span of all but the last vector is
complete by induction, hence closed. The last vector has positive
distance from this span, bounding its coefficient in any combination.
The ultrametric inequality and induction then bound all other
coefficients. The reverse bound is the triangle inequality. Consequently
finite-dimensional normed spaces are complete and their coordinate maps
bounded. Existence of an absolute value on a finite field extension can
also be seen by giving the finite algebra its finite-module Banach
topology, using CLD 11.4.4, and taking a point of its nonempty affinoid
spectrum. Its kernel is zero since the algebra is a field. Uniqueness
follows from norm equivalence: compare two extending norms on \(z^m\)
and take \(m\)th roots. Hence all finite field automorphisms preserve
the absolute value, and every finite field here is complete.

For a \(K\)-affinoid algebra \(A\), tensor the bounded coordinate maps
of \(E\) with \(A\). The tensor norm upper bound and these coordinate
maps identify its topology with a finite sum of copies of \(A\). Already
complete, this proves \[
             B=A\widehat\otimes_K E=A\otimes_K E.        \tag{5}
\] This is the actual algebra, including nilpotents. Taking the kernel
of \(E\to\bigoplus_{g\in G}E\), \(z\mapsto(gz-z)_g\), and tensoring over
the field gives \(B^G=A\). Invariant units descend as units since their
inverses are invariant too. There is no field averaging or division by
\(d\). An affinoid element nonzero at every analytic point is a unit:
otherwise a maximal ideal containing it and spectral surjectivity give a
zero.

The fibre over \(x\) is \(\mathcal M(\mathcal H(x)\otimes_K E)\). A
primitive separable polynomial factors without repeated factors over
\(\mathcal H(x)\), so the tensor is a finite product of fields. Each has
one analytic point by the existence and uniqueness argument above. The
\(G\) action on these finitely many points is transitive: two orbits
would supply a nontrivial invariant idempotent, but the invariant
algebra is \(\mathcal H(x)\). Stabilizers can be nontrivial; no local
covering-space description is claimed or used.

Inverse images of affinoid ordinary interiors give goodness upstairs.
Strict presentations base change to strict presentations. Base change of
the closed diagonal gives separatedness. Boundarylessness follows from
Section 2(2). In particular all carriers are Hausdorff. The inverse
image of any compact \(F\subset X\) is compact: cover \(F\) by finitely
many interiors of affinoid neighbourhoods \(V_i\); its inverse image is
closed in the finite union of the compact \(V_{i,E}\).

\textbf{Full-fibre shrinking.} In a downstairs strict affinoid
neighbourhood \(V_0\) of \(x\), let \(N\subset V_{0,E}\) be a relative
ordinary open containing the whole fibre. The compact complement has
closed image excluding \(x\). Choose a smaller strict affinoid
neighbourhood \(V\) of \(x\) in the complement of that image and inside
\(\operatorname{int}_X(V_0)\). Then \(V_E\subset N\). Thus finitely many
unit or strict norm conditions holding on the whole fibre persist on one
neighbourhood. This includes conditions on elements of zero seminorm at
\(x\) and never replaces \(A\) by its reduction.

For completeness, purity is checked at minimal components. In a strict
affinoid \(\mathcal M(A)\) in \(X\), each minimal-prime component has
dimension one: choose a point with kernel that prime by spectral
surjectivity and apply the local dimension formula. Let \(J\) be the
nilradical. It is nilpotent since \(A\) is noetherian. The injection
\(A/J\hookrightarrow\prod_P\operatorname{Frac}(A/P)\) stays injective
after tensoring with a separable field. Each field tensor is a product
of fields, hence \((A/J)\otimes E\) is reduced and the nilradical of
\(B\) is exactly \(J\otimes E\). The finitely many kernels of its maps
to those field factors have zero intersection. Their product is zero, so
a minimal prime contains one of these kernels, and by minimality equals
it. Its contraction is the corresponding minimal \(P\) of \(A\). The
finite integral injection \(A/P\hookrightarrow B/Q\) preserves Krull
dimension. One elementary justification is lying over and going up from
the determinant argument for a finite module over a local ring:
\(\mathfrak pD_{\mathfrak p}\ne D_{\mathfrak p}\), since otherwise the
determinant of \(I-M\), \(M\) with entries in \(\mathfrak p\), is a unit
annihilating that nonzero module. Choose a maximal ideal above it and
then repeat after quotienting by a chosen prime. For incomparability
quotient by the smaller prime and invert the nonzero elements of the
base domain. The resulting finite integral domain over a field is a
field, since a least-degree monic equation for a nonzero element
expresses its inverse. A strictly larger prime of the same contraction
could not survive, nor could a nonzero element of a domain vanish under
that localization. Going up lifts chains and incomparability makes
contractions strict. The dimensions agree. Every minimal component
upstairs therefore has dimension one, proving purity. The same proof
works for any finite separable stage. Reduction here is a dimension
check only; all measures below use \(B\) itself.

Finally consider the same carrier \(Z=X_E\) with its restricted
\(K\)-analytic structure \(Z_0\). A finite \(E/K\) is \(K\)-affinoid; an
\(E\)-affinoid presentation is therefore \(K\)-affinoid under finite
field restriction. The finite-module topology is unchanged. Rational
domains are defined by the same elements and inequalities and have the
same analytic rings. They give a common local basis and the same unit
germs and compact moment images. The moment nullity criterion of Section
2(3) identifies the ordinary form sheaves, their exterior calculus and
supports. Each \(g\in G\) is a bounded \(K\)-algebra automorphism of
(5), gluing to a \(K\)-analytic automorphism of \(Z_0\). Its isomorphism
pushforward formula, followed by cancellation of the same-carrier factor
\(d\), gives for a smooth compact top form \(\omega\) \[
             g_*((g^*\omega)_{\mathrm{Rad},E})
                            =\omega_{\mathrm{Rad},E}.    \tag{6}
\] This is different from (4). Both are valid with actual generic
lengths.

\subsection{4. Compact partitions and full-fibre scalar
descriptions}\label{compact-partitions-and-full-fibre-scalar-descriptions}

Let compact \(F\) have an ordinary open cover. At every point of \(F\),
apply Section 2(4)'s bump theorem to the singleton in a chosen cover
member. Squaring the bump makes it nonnegative, still equal to one near
that point and compactly supported there. Select finitely many whose
positive loci cover \(F\), and group them by cover index as \(\psi_i\).
Then \(s=\sum_i\psi_i>0\) near \(F\). Take a bump \(\chi=1\) near \(F\)
compactly supported in \(\{s>0\}\) and put \[
                \theta_i=\chi\psi_i/s                  \tag{7}
\] there, extending by zero. Products, smooth composition, and
reciprocal on \(s>0\) are ordinary smooth operations in moment
coordinates. The compact support of \(\chi\) lies inside \(\{s>0\}\), so
extension by zero is smooth also at its complement. Each support is
compact and contained in its cover member. The sum is one near \(F\).
This proves the needed finite compact partition without any noncompact
fixed-cover-index assertion.

Let \(\beta\) be an invariant ordinary smooth scalar function near a
whole fibre over \(x\). Its fibre has finitely many points. Choose
pairwise disjoint ordinary neighbourhoods \(N_i\) of them with local
representations \(\beta=\varphi_i(t_{f_{i,1}},\ldots)\), using Proposition
15.1.1(b) in scalar bidegree. Shrink \(N_i\) so the coordinate values are in a relatively
compact Euclidean region on which \(\varphi_i\) is smooth. Full-fibre
shrinking gives a strict affinoid \(V\) with \(V_E\subset\bigcup_iN_i\).
The pieces \(C_i=V_E\cap N_i\) are disjoint ordinary opens and closed in
\(V_E\). They are compact. Analytic gluing gives idempotents \(e_i\)
which are one on \(C_i\) and zero on the other pieces, and gives a
global unit \(\widetilde f_{i,k}\) equal to \(f_{i,k}\) on \(C_i\) and
one elsewhere. Its inverse glues in the same way. These are elements of
the actual affinoid algebra \(B\), not formal topological selectors.

Choose \(a\in K^\times\) with \(|a|\ne1\). The analytic unit
\(s_i=e_i+a(1-e_i)\) has \(t_{s_i}=0\) on \(C_i\) and
\(t_{s_i}=-\log|a|\) elsewhere. Choose a real smooth \(q\) equal to one
near zero and zero near \(-\log|a|\). Extend \(\varphi_i\) by a
Euclidean cutoff equal to one near its compact coordinate image. Thus \[
 \beta=\sum_i q(t_{s_i})\varphi_i(t_{\widetilde f_{i,1}},\ldots)
                          =\Phi(T)\quad\text{on }V_E,   \tag{8}
\] for one global unit tuple \(T\) and a Euclidean smooth \(\Phi\). All
group translates of these units exist on \(V_E\); adjoining them gives
one fixed tuple. The real average \(|G|^{-1}\sum_g\Phi(T\circ g)\)
equals \(\beta\) there. This construction works for every orbit and
stabilizer, including singleton fibres. It does not assume exact
invariant ordinary scalar descent.

\subsection{5. Polarization removes the orbit-correlation
difficulty}\label{polarization-removes-the-orbit-correlation-difficulty}

For a smooth Euclidean function on a neighbourhood of a compact box,
there are real polynomials converging through order two on that box.
Here is a proof of precisely the version used. Extend it by a compact
Euclidean cutoff, convolve with a Gaussian, and differentiate under the
integral. Uniform continuity of the derivatives through order two and
the Gaussian tail give uniform derivative convergence on the box. For
fixed convolution width, the Gaussian and its first two derivatives are
uniform limits of their Taylor polynomials on the compact difference of
the box and the original coefficient support. Integrate these
polynomials to get polynomial approximants through order two, then
choose widths and truncation degrees diagonally.

For \(k\ge1\) and real symbols \(z_1,\ldots,z_k\), including repeated
symbols, inclusion-exclusion and the multinomial expansion give \[
 z_1\cdots z_k=\frac1{k!}\sum_{J\subseteq\{1,\ldots,k\}}
            (-1)^{k-|J|}\left(\sum_{j\in J}z_j\right)^k. \tag{9}
\] Indeed each term missing one symbol cancels, and the only surviving
degree-\(k\) term containing every symbol is their product with
coefficient \(k!\). For unit logarithms, a sum is the logarithm of the
corresponding product unit. Consequently the real group average of
\textbf{every} scalar polynomial in an arbitrary tuple of unit
logarithms is a finite real linear combination of constants and terms \[
        T_{f,k}=\frac1d\sum_{g\in G}(t_{gf})^k,          \tag{10}
\] with actual product units \(f\). This recovers mixed orbit
correlations; no claim that separate symmetric sums for the original
individual units suffice is used. The number of product units may grow
with polynomial degree. Section 9 handles this explicitly; no common
neighbourhood for infinitely many product units is assumed.

\subsection{6. Actual root valuations and all coefficient-zero
cases}\label{actual-root-valuations-and-all-coefficient-zero-cases}

For a unit \(f\in B^\times\) set \[
 P_f(Z)=\prod_{g\in G}(Z-gf)
             =Z^d+\sum_{i=0}^{d-1}c_iZ^i.               \tag{11}
\] Its coefficients lie in \(B^G=A\); \(c_0\) is a unit. At an analytic
point \(y\), the roots \(gf(y)\) are actual nonzero elements of
\(\mathcal H(y)\) with multiplicities indexed by \textbf{all} group
elements. This includes repeated roots and arbitrary stabilizers.
Neither a separable orbit polynomial nor characteristic-zero
coefficients are required.

We prove the root-valuation rule rather than importing Newton-polygon
recognition. For a polynomial \(\prod_{j=1}^d(Z-r_j)\) over any valued
field with nonzero roots, sort \(u_j=-\log|r_j|\) in nondecreasing
order. The coefficient of \(Z^{d-k}\) is a sum of products of \(k\)
roots. Every summand has valuation at least \(u_1+\cdots+u_k\), so its
coefficient has valuation at least that number, allowing infinity. If
\(u_k<u_{k+1}\), the subset of the \(k\) roots of smallest valuation is
unique and its product has strictly greater norm than every other
summand. Ultrametricity makes the coefficient valuation exactly that
sum. Equality also holds at \(k=0,d\). These equality points are
precisely the ends of constant-valuation blocks. The convex polygon
joining them lies below all the coefficient points and has its
successive slopes equal to the root valuations in reversed sign when
indexed by coefficient exponent. Thus the lower convex hull of \[
          (i,v(c_i))\ (0\le i\le d),\qquad v(c_d)=0      \tag{12}
\] has each segment \(i<j\) contributing \(j-i\) roots of valuation
\((v(c_i)-v(c_j))/(j-i)\). Repeated roots, ties, and every cancellation
are covered by the strict-gap argument. No division occurs in the valued
field.

Near a fixed downstairs \(x\), choose \(M\) so all \(|t_{gf}|<M\) on its
finite fibre and retain this inequality by full-fibre shrinking. Then
\(|t_{c_0}|<dM\). If \(c_i(x)=0\) for \(0<i<d\), shrink further until \[
                      |c_i|<\exp(-dM-1).                \tag{13}
\] Its valuation lies strictly above the chord from \((0,t_{c_0})\) to
\((d,0)\), since that chord's height is at most \(dM\). Hence it cannot
belong to the lower hull and may be omitted. If \(c_i(x)\ne0\), shrink
to make it a downstairs unit: choose a strict affinoid neighbourhood
inside its nonzero locus, then use the unit criterion of Section 3. Keep
\(c_0,c_d=1\) in all cases. The retained index set
\(I\subseteq\{0,\ldots,d\}\) therefore includes both ends. Its other
coefficients are actual units; all omitted coefficients obey (13)
throughout this neighbourhood. Nilpotent coefficients are included in
the zero-seminorm case. No coefficient with value zero is logarithmed,
and small coefficients are not simply assumed harmless.

\subsection{\texorpdfstring{7. A \(C^1\) piecewise polynomial and its
tangential
mollification}{7. A C\^{}1 piecewise polynomial and its tangential mollification}}\label{a-c1-piecewise-polynomial-and-its-tangential-mollification}

For arbitrary real heights \(z_i\), \(i\in I\), with \(z_d=0\), form the
lower hull of the finitely many points \((i,z_i)\). If consecutive
vertices are \(i<j\), put \(m=j-i\) and \(u=(z_i-z_j)/m\). Define \[
               H_{I,k}(z)=\sum_{\text{hull segments }i<j}
                                 m\,u^k.                \tag{14}
\] Collinear vertices may be retained or omitted: (14) is unchanged.
There are finitely many possible hull chains. Each is selected by linear
weak inequalities in the heights, so (14) is a polynomial on each region
of a finite polyhedral subdivision of the whole height space. Equations
(11)--(13) prove at \textbf{every} point upstairs that \[
               T_{f,k}=d^{-1}H_{I,k}(t_{c_i}).            \tag{15}
\]

We need more than continuity of this formula. For a general smooth
one-variable \(\phi\), the segment term \(m\phi((z_i-z_j)/m)\) has
gradient \(\phi'(u)(e_i-e_j)\) and Hessian \[
            \frac{\phi''(u)}m(e_i-e_j)\otimes(e_i-e_j).  \tag{16}
\] If a vertex is inserted or deleted on a collinear segment, all the
affected \(u\) equal one common value. The gradients then telescope:
\(\sum\phi'(u)(e_i-e_j)\) is exactly the gradient of the single merged
segment. This remains true if several vertices are collinear. At a
height configuration on a chamber boundary, every limiting hull chain
refines the unique lower hull by such collinear vertices. Values and
gradients therefore agree for every limiting chamber. A finite piecewise
smooth function with these matching values and gradients is \(C^1\):
integrate the matching, continuous candidate gradient on short line
segments subdivided by the finitely many chamber faces; segments
contained in faces follow by a small translation and continuity. The
usual differentiability remainder follows from continuity of that
gradient. Thus (14) is \(C^1\).

On any compact neighbourhood, slopes \(u\) are bounded. Each Hessian
(16) is bounded by \(2\sup|\phi''|/m\) in operator norm; there are at
most \(d\) segments. For \(\phi(u)=u^k\), there is a finite bound for
all chamber Hessians on that neighbourhood. The gradient is locally
Lipschitz by integrating these Hessians along line segments, using its
continuity at faces as above. The weak Hessian is the bounded piecewise
Hessian, with no singular face measure: integrate by parts on each
polyhedral region; the boundary contributions from the continuous
gradient cancel in pairs. Lower-dimensional intersections give no
additional boundary term. This proves all assertions used in
mollification; bounded third or fourth derivatives are unnecessary.

Choose a nonnegative even \(C^\infty\) kernel \(\rho\) supported in the
unit ball with integral one. Let \(H_\varepsilon=H*\rho_\varepsilon\).
The local gradient Lipschitz bound yields uniformly on fixed compacts
\(H_\varepsilon\to H\) and \(\nabla H_\varepsilon\to\nabla H\); in fact
the errors are \(O(\varepsilon^2)\) and \(O(\varepsilon)\). The first
bound follows from the first-order remainder and the kernel's zero first
moment. Its Hessian is the convolution of the weak Hessian and is
uniformly bounded on those compacts.

\textbf{Tangential second derivative lemma.} Let an open segment \(L\)
lie in the relative interior of a face \(F\) of the subdivision of
\(H\), with tangent vector \(v\) in the affine direction of \(F\). At
every point \(z\in L\), \[
        \nabla^2H_\varepsilon(z)(v,v)
                  \longrightarrow (H|_L)''(z;v,v).       \tag{17}
\] To prove this, only chambers incident to \(F\) meet a sufficiently
small ball about \(z\). Let their polynomial formulas be \(Q_a\). Each
agrees with \(H\) on a relatively open neighbourhood in \(F\), so
\(\nabla^2Q_a(z)(v,v)\) is the same tangential derivative of that
restriction. For points \(w\) tending to \(z\) through any such chamber,
the corresponding tangential Hessian tends to this common value.
Convolving the bounded, piecewise Hessian proves (17). If the line is
mapped to one point in height space, the tangent is zero and the
statement is immediate. This proof deals with segments \textbf{entirely
contained in collision walls}; they are not dismissed as measure zero. A
compact moment segment maps affinely to height space and can be
subdivided at finitely many face preimages, so the conclusion holds
almost everywhere on every such segment, including those of deficient
height-map rank.

\subsection{8. Calibrated differentiated approximation of a fixed
trace}\label{calibrated-differentiated-approximation-of-a-fixed-trace}

Fix \(f,k\) and one neighbourhood on which (13)--(15) hold. The
functions \[
                b_\varepsilon=d^{-1}H_\varepsilon(t_{c_i})
                                                               \tag{18}
\] are ordinary smooth downstairs, using actual unit moments. On a
compact upstairs domain add to one moment tuple the units \(gf\), the
pullbacks of all retained coefficients, and any fixed cutoff's units.
Its image is a finite polyhedral set. Refine its one-dimensional cells
by the height-fan preimages and by a common skeletal calibration
subdivision. On each segment, (15) is an identity of functions of affine
coordinates. Its right side restricts to a polynomial on the segment,
whose first and second tangent derivatives are therefore those of the
ordinary upstairs polynomial \(T_{f,k}\).

The errors \(e_\varepsilon=T_{f,k}-\pi^*b_\varepsilon\) and their first
derivatives converge uniformly tangentially. Their second tangential
derivatives are uniformly bounded and converge almost everywhere on
every segment by (17). No differentiation of (15) as a nonsmooth
ordinary section is performed; the derivatives are compared on actual
moment cell tangents.

Here is why this is a total-variation statement on the actual carrier.
For the enlarged tuple \(t_1,\ldots,t_r\) put
\(\zeta=\sum_a d't_a\wedge d''t_a\), with the positive top-order
convention. On a segment with tangent \(v\), a coefficient matrix
\(A=(A_{ab})\) contracts to \(\sum A_{ab}v_av_b\) times its positive
calibration; \(\zeta\) contracts to \(\sum v_a^2\) times the same
calibration. Hence \(|A_{ab}|\le M\) gives \[
       |\eta_{\mathrm{Rad}}|\le rM\zeta_{\mathrm{Rad}},    \tag{19}
\] by \(\big(\sum|v_a|\big)^2\le r\sum v_a^2\). Equivalently use bounded
tangential coefficients in unit tangent frames on the finitely many
segments. The estimate is made on the supporting skeletal cells before
pushforward. Section 2(5) gives finite total mass of \(\zeta\) on the
compact domain and a constant finite positive density on each moment
segment. Thus bounded tangential error converging almost everywhere has
carrier total variation tending to zero by bounded convergence on those
finite calibrated segments. Zero-dimensional faces carry no top measure.
Multiplicities and generic lengths multiply both sides of (19); none is
removed or replaced by a point count.

For any \textbf{fixed} ordinary smooth cutoff \(\lambda\) supported in
the neighbourhood, the product rule for \(D(\lambda e_\varepsilon)\)
uses the scalar error, its two first derivatives, and its mixed second
derivative, together with fixed bounded derivatives of \(\lambda\). The
preceding convergence and (19) show \[
       \big\|(D((\lambda\circ\pi)e_\varepsilon))_{
                          \mathrm{Rad},E}\big\|_{\mathrm{TV}}
                                 \longrightarrow0.       \tag{20}
\] Cover a fixed compact cutoff support by finitely many compact moment
neighbourhoods to use this argument; their enlarged tuples and
calibration refinements are fixed during this limit. Signed fibre
cancellation is not used to infer (20).

For a fixed finite combination of traces (10), Sections 6--8 apply near
every downstairs point, with finitely many units, omitted coefficients
and strict inequalities at that point. If a fixed downstairs cutoff
\(\chi\) has compact support in the original unit domain, cover its
support by finitely many of these local neighbourhood interiors, and
choose an inner compact partition \(\theta_a\) by (7). Let
\(b_{a,\varepsilon}\) be the local combination of (18). Put \[
        a_\varepsilon=
              \sum_a\widetilde{\chi\theta_a b_{a,\varepsilon}},     \tag{21}
\] where tilde means extension by zero. This is ordinary smooth and
supported in \(\operatorname{supp}\chi\). A product supported compactly
inside its ordinary open is identically zero on a neighbourhood of each
point outside that open, so this extension is smooth and has no boundary
derivative term. The sum-one identity near \(\operatorname{supp}\chi\)
and (20), for this finite fixed cover, prove \[
 \big\|(D((\chi\circ\pi)T-\pi^*a_\varepsilon))_{
                      \mathrm{Rad},E}\big\|_{\mathrm{TV}}\to0       \tag{22}
\] for the given invariant scalar polynomial \(T\). All supports here
lie in \(\pi^{-1}(\operatorname{supp}\chi)\), independently of
\(\varepsilon\). There is no assertion of positive-test preservation.

\subsection{9. Every compact ordinary scalar test; the two limits and
support}\label{every-compact-ordinary-scalar-test-the-two-limits-and-support}

Let invariant \(\beta\in A_c^{0,0}(X_E)\) be arbitrary, and set
\(F=\pi(\operatorname{supp}\beta)\). It is compact. Cover \(F\) by
finitely many ordinary interiors \(U_i\) of strict affinoids \(V_i\) on
which (8) holds, and choose outer cutoffs \(\chi_i\) from (7), compactly
supported in \(U_i\), with sum one near \(F\). Put
\(C=\bigcup_i\operatorname{supp}\chi_i\). It is a fixed compact set
containing \(F\).

On \(V_{i,E}\) approximate the fixed Euclidean coefficient \(\Phi_i\)
through order two on a box containing the fixed unit moment image by
polynomials \(P_{i,j}\). Choose the box to contain all translated tuple
images as well. Let \(T_{i,j}\) be the real group average of
\(P_{i,j}(T_i)\). Since \(\beta\) is invariant, it is also the average
of \(\Phi_i(T_i)\), so \(\beta-T_{i,j}\) tends to zero in coefficient
\(C^2\) on a \textbf{fixed upstairs} moment tuple containing all the
translates. Include fixed moment descriptions of \(\chi_i\circ\pi\) on a
finite compact cover of its support. The product rule and (19) give \[
 \big\|(D((\chi_i\circ\pi)(\beta-T_{i,j})))_{
                     \mathrm{Rad},E}\big\|_{\mathrm{TV}}\to0.       \tag{23}
\]

For each fixed \(j\), use (9) to express \(T_{i,j}\) as finitely many
traces. Apply (21)--(22) with outer cutoff \(\chi_i\) to obtain
\(a_{i,j}\) supported in \(\operatorname{supp}\chi_i\) and with \[
 \big\|(D((\chi_i\circ\pi)T_{i,j}-\pi^*a_{i,j}))_{
                     \mathrm{Rad},E}\big\|_{\mathrm{TV}}<1/j       \tag{24}
\] (or \(1/(ij)\) after any fixed finite enumeration). Each
fixed-polynomial local cover, inner partition and moment refinement is
chosen \textbf{before} its smoothing limit. It may depend on \(j\). Its
derivative bounds may grow with \(j\); smoothing is chosen after those
bounds. This causes no problem because (23) is evaluated with the fixed
outer cutoffs, before the inner localization of \(T_{i,j}\). The varying
inner cutoffs never multiply the polynomial approximation error in (23),
and every approximate support remains inside the same outer \(C\).
Neither a degree-uniform estimate nor an infinite simultaneous unit
shrink is asserted.

Set \(\alpha_j=\sum_i a_{i,j}\). Globally
\(\beta=\sum_i(\chi_i\circ\pi)\beta\), since the sum is one near its
support and it vanishes off that support. Equations (23)--(24) and the
finite triangle inequality prove (2). The support of every
differentiated error is contained in \(\pi^{-1}(C)\), compact by Section
3. Thus for \textbf{every} continuous finite multiplier \(h\), \[
       B_{h\circ\pi}(\beta-\pi^*\alpha_j)\longrightarrow0,
                                                               \tag{25}
\] using its boundedness on that compact set. This proves the actual
ordinary scalar test gate for curves, not merely \(C^0\) density of
invariant functions.

\subsection{10. Forward Galois ascent and actual normal-closure
descent}\label{forward-galois-ascent-and-actual-normal-closure-descent}

For a general compact ordinary scalar test \(\beta\) upstairs, (6),
ordinary exterior functoriality on the restricted carrier, and
invariance of \(h\circ\pi\) show
\(B_{h\circ\pi}(g^*\beta)=B_{h\circ\pi}(\beta)\). Its real average
\(\beta^{\mathrm{av}}=d^{-1}\sum_g g^*\beta\) has the same pairing. Its
support is in a compact finite group saturation. Real averaging does not
divide by \(d\) in \(K\), including when \(\operatorname{char}K\)
divides \(d\).

For every downstairs compact test \(\alpha\), its pullback is compact.
The scalar projection formula (4), with \textbf{no degree factor}, gives
\[
  B_{h\circ\pi}(\pi^*\alpha)
     =\int_{X_E}(h\circ\pi)(\pi^*D\alpha)_{\mathrm{Rad},E}
     =\int_X h(D\alpha)_{\mathrm{Rad},K}=B_h(\alpha).     \tag{26}
\] Smooth exterior functoriality suffices; no continuous-current
pullback has been presumed. If the downstairs current is zero,
approximate \(\beta^{\mathrm{av}}\) by (2). Its approximate pairings
vanish by (26), and (25) proves that its pairing, hence that of every
upstairs test, vanishes. This proves forward ascent for every finite
Galois extension, including wild ones, on all curves in the stated
scope.

Conversely, (26) alone proves reverse descent whenever the upstairs
current is zero. This calculation works for every finite separable
extension, without a group, by compactness, purity, and (4). It is an
actual proof of the descent needed below, not a use of an earlier note.

For arbitrary finite separable \(E/K\), take the finite Galois normal
closure \(N/K\) constructed in Section 3, with its unique extending
absolute value and completeness. Forward Galois ascent gives zero
current on \(X_N\). The space \(X_E\) retains every ambient hypothesis
and pure dimension one by Section 3, and \(N/E\) is finite separable.
There is a canonical analytic identification \((X_E)_N=X_N\): on
affinoids it is the associative tensor isomorphism
\((A\otimes_K E)\otimes_E N=A\otimes_K N\) with the equivalent
finite-module topologies. The pullback of \(h\circ\pi_E\) is the
original pullback on \(X_N\). Apply the just-proved reverse calculation
to \(N/E\), testing each compact ordinary function on \(X_E\), to obtain
zero current on \(X_E\). This proves (1) at every finite separable
stage. It is the finite dependency, before the completed stage in
Section 13. Ordinary opens retain the stated assumptions and the same
argument applies to their scalar inverse images.

\subsection{11. Finite inseparability and arbitrary finite algebraic
stages}\label{finite-inseparability-and-arbitrary-finite-algebraic-stages}

\subsubsection{11.1. Frobenius, topology and exact scalar
tests}\label{frobenius-topology-and-exact-scalar-tests}

Let \(k\) be complete, nontrivially valued, of characteristic \(p>0\),
and let \(F/k\) be finite purely inseparable. Choose \(q=p^e\) with
\(F^q\subset k\). Section 3's bounded-coordinate argument applies to
\textbf{every} finite extension, and gives
\(A\widehat\otimes_k F=A\otimes_k F\) with its finite-module topology.
Consequently every \(u=\sum a_i\otimes\lambda_i\) satisfies \[
                    u^q=\sum a_i^q\lambda_i^q\in A.       \tag{27}
\] This is an identity in the actual algebra, including nilpotents.

Scalar projection is surjective by the affine foundation in Section 2;
Section 12.1 also gives a direct proof for all fields used here. If two
points of \(\mathcal M(A\otimes F)\) project to \(x\), (27) forces their
values on every \(u\) to equal \(|u^q(x)|^{1/q}\). They are the same
seminorm. The projection is thus a continuous bijection from a compact
space to a Hausdorff space, hence a homeomorphism. On a good separated
space these affine homeomorphisms glue on the ordinary interiors of
downstairs affinoid neighbourhoods. No covering-space assertion is used.

For a unit \(f\in(A\otimes F)^\times\), (27) gives \(f^q=g\in A\). The
function \(g\) has no analytic zero downstairs, by surjectivity, so it
is an actual unit by Section 3's unit criterion. On the upstairs carrier
\(t_f=q^{-1}t_g\). Now let \(\beta\in A_c^{0,0}(X_F)\) be arbitrary. Its
support has compact image \(C\) under the homeomorphism. For each point
of \(C\), select an ordinary local description of \(\beta\), then a
downstairs affinoid neighbourhood \(V\) whose full inverse image lies in
that description open. The inverse image is \(V_F\), so all coordinate
units restrict to elements of its actual affinoid algebra. Taking their
\(q\)th powers gives an ordinary downstairs scalar description
\(\varphi(q^{-1}t_{g_1},\ldots,q^{-1}t_{g_m})\) on
\(\operatorname{int}_X V\). Equality of numerical smooth scalar
functions is equality of their ordinary \((0,0)\) sections, as explained
in Section 12.3. Cover \(C\) by finitely many such opens, apply the
compact partition of Section 4, multiply these local descriptions by its
cutoffs, and extend by zero. This produces \(\alpha\in A_c^{0,0}(X)\)
with \(\beta=\pi_F^*\alpha\) everywhere. Compact support inside each
open makes the extension ordinary smooth; no derivative is extended
through an exposed boundary.

Goodness, strictness, separatedness and boundarylessness are preserved
as in Section 3. Purity is also preserved, and here is a proof that does
not assume equality of total dimensions gives local purity. Given an
upstairs point \(y\), a downstairs affinoid neighbourhood of its image
has dimension one, and its scalar extension has dimension one, giving
the upper local bound. If local dimension at \(y\) were zero, take this
affinoid upstairs and remove all its positive-dimensional irreducible
components, none of which contains \(y\). An affinoid neighbourhood
inside the remaining ordinary open has dimension zero by CLD 11.2.3's
component/domain formula. By the homeomorphism choose a downstairs
affinoid neighbourhood whose full inverse image is inside that open. Its
downstairs dimension is one (it is a domain in the pure curve), and its
upstairs total dimension is one by scalar dimension invariance, a
contradiction. Existence of the smaller downstairs affinoid
neighbourhood uses goodness. This argument only checks dimension; it
never reduces the measure.

On these pure curves \(D\beta=\pi_F^*D\alpha\), and the smooth
top-measure identity (4) proves \[
             B_{h\circ\pi_F}(\beta)=B_h(\alpha).         \tag{28}
\] All supports are compact. The converse is obtained by taking
\(\beta=\pi_F^*\alpha\) for an arbitrary compact downstairs scalar test.
Thus finite purely inseparable zero-Hessian ascent and descent are
proved here directly. No earlier inseparable manuscript is a premise.

\subsubsection{11.2. The finite mixed-field
decomposition}\label{the-finite-mixed-field-decomposition}

Let \(E/K\) be any finite algebraic extension. In characteristic zero
every irreducible polynomial has nonzero derivative, and is separable,
so Sections 3--10 apply. In characteristic \(p\), choose finitely many
generators \(a_1,\ldots,a_r\) of \(E/K\). For each irreducible minimal
polynomial, repeatedly remove common factors \(p\) from its exponents
until it has the form \[
                f_i(Z)=g_i(Z^{p^{e_i}}),\qquad g_i'\ne0.
                                                               \tag{29}
\] The polynomial \(g_i\) is irreducible: a factorization of \(g_i\)
would give a factorization of \(f_i\) by substitution. Since \(g_i\) is
irreducible and its derivative is nonzero of smaller degree,
\(\gcd(g_i,g_i')=1\); its roots are distinct. Therefore
\(b_i=a_i^{p^{e_i}}\) is separable over \(K\). The compositum
\(S=K(b_1,\ldots,b_r)\subset E\) is finite separable: adjoining a root
of a separable polynomial remains separable over an intermediate field
because its minimal polynomial divides that polynomial; extending
embeddings successively gives degree-many embeddings of the tower. An
intermediate simple field has degree-many embeddings as well: each
embedding has at most the remaining extension degree lifts, while the
entire tower has its full degree of embeddings. Thus its minimal
polynomial is separable. This proves the field facts used here without
assuming separability of an arbitrary algebraic stage.

With \(e=\max e_i\), every \(a_i^{p^e}\in S\), and taking \(p^e\)th
powers in a rational expression in the \(a_i\) shows
\(E^{p^e}\subset S\). Hence \(E/S\) is finite purely inseparable. Both
fields are complete under the prescribed absolute value by the finite
norm proof in Section 3. This gives the \textbf{actual embedded} tower
\(K\subset S\subset E\), not an abstract factorization of degrees.

The finite tensor associativity with bounded inverse maps, proved in
Section 12.1, identifies \((X_S)_E=X_E\) as analytic spaces, with the
actual algebra and not its reduction. \(X_S\) has all ambient hypotheses
and pure curve dimension by Section 3. Finite separable ascent to \(S\),
then the proved Frobenius ascent to \(E\), proves (1) at every finite
algebraic stage. Reverse descent uses (4) directly. Generic lengths are
those of the actual tensors throughout. No degree factor is introduced
into scalar projection.

\subsection{12. Completion, rational neighbourhoods and every compact
scalar
test}\label{completion-rational-neighbourhoods-and-every-compact-scalar-test}

\subsubsection{12.1. Fields, completed tensors, surjectivity and
compactness}\label{fields-completed-tensors-surjectivity-and-compactness}

Write \(M=\bigcup_E E\), where \(E\) ranges over its finite
subextensions of \(K\). It is a directed union: a finite list of
algebraic elements generates a finite field, by successive
multiplication of finite degrees. Each \(E\) is complete by Section 3.
We never regard \(M\) as a complete analytic field unless it happens
already to be complete. Its completion \(L\) is a complete valued field.
For a \(K\)-affinoid algebra \(A\), the image of \(A\otimes_K M\) is
dense in \(A_L=A\widehat\otimes_K L\). Indeed finite algebraic tensors
are dense by definition of completion; a tensor
\(\sum a_j\otimes\lambda_j\) can be approximated by replacing its
finitely many \(\lambda_j\in L\) with \(m_j\in M\), with tensor error at
most \(\max_j\|a_j\||\lambda_j-m_j|\). Approximation statements refer to
the \textbf{image}, so no injectivity or isometry of an uncompleted
infinite tensor is presumed.

Here is the associative tensor argument actually required. For finite
\(E/K\) with a basis \(e_i\), the bounded coordinate maps give
\(A_E=A\otimes_K E\) with its finite-module topology. The maps on
algebraic tensors \[
 (a\otimes e)\otimes\lambda\mapsto a\otimes e\lambda,
 \qquad a\otimes\lambda\mapsto(a\otimes1)\otimes\lambda
                                                               \tag{30}
\] extend continuously to inverse maps
\((A_E)\widehat\otimes_E L\simeq A\widehat\otimes_K L\). For the first
map, the bounded multiplication map
\((b,\lambda)\mapsto\sum_i a_i\otimes e_i\lambda\), where
\(b=\sum_i a_i\otimes e_i\), is bounded in \(b\) by the coordinate
estimates and in \(\lambda\) by the tensor inequality. It is balanced
over \(E\). The reverse map is bounded by the tensor inequality and the
bounded inclusion \(A\to A_E\). These maps are inverse on dense
elementary tensors, hence inverse on the completions. The same proof
works when \(A\) is affinoid with arbitrary radii, and for two finite
stages. Thus all stage projections and inverse images of rational
domains used below are actual analytic base changes. No associativity
claim for an infinite uncompleted Banach tensor is needed.

For clarity, scalar surjectivity in the present completed-algebraic
scope also has a direct seminorm proof. Given \(x\in\mathcal M(A)\), put
\(H=\mathcal H(x)\). Embed the algebraic extension \(M\) into an
algebraic closure \(\overline H\) over \(K\). Such an embedding is
obtained by successively adjoining roots, or by maximizing partial
embeddings and adjoining a root of the minimal polynomial if the domain
is not all of \(M\). Equip \(\overline H\) with the extending absolute
value: the finite extension existence and uniqueness proof of Section 3
makes the norms on its finite \(H\)-subfields compatible. On every
finite \(K\)-subfield of \(M\) this value is the specified one, again by
uniqueness. Completing \(\overline H\) gives a complete field \(F\) into
which the isometric embedding \(M\to F\) extends to \(L\to F\). For
\textbf{every} tensor representation with \(\lambda_j\in L\), evaluation
satisfies \[
 \left|\sum_j a_j(x)\lambda_j\right|
       \le \max_j\|a_j\||\lambda_j|.
                                                               \tag{31}
\] Taking the infimum over all representations bounds it by the full
tensor seminorm. It extends to \(A_L\) and gives an analytic point over
\(x\). This proof handles the separated quotient in the tensor
completion and uses no claimed isometry of \(A\otimes M\). The identical
proof over \(E\) gives surjectivity \(X_L\to X_E\) and surjectivity on
every scalar-extended affinoid domain. On good spaces these affine facts
cover the whole space.

The inverse image of a compact \(C\subset X_E\) is compact: finitely
many ordinary interiors of affinoid neighbourhoods cover \(C\). Its
inverse image is closed in the finite union of their compact scalar
extensions, since the spaces are Hausdorff. This works for the infinite
completed projection. It does not assert a finite fibre, a finite Galois
cover or any infinite group averaging.

For a nonzero \(\lambda\in L\), choose \(m\in M\) with
\(|m-\lambda|<|\lambda|\). Then \(|m|=|\lambda|\). Thus
\(|L^\times|=\bigcup_E|E^\times|\). This observation handles any
value-group radius if desired, but the rational argument below allows
\textbf{all positive real radii}, so it needs no strictifying extension
and no approximation of a radius.

\subsubsection{12.2. Exact rational and strict-locus
descent}\label{exact-rational-and-strict-locus-descent}

Let \(A\) be affinoid over a finite stage \(E_0\), set
\(B=A\widehat\otimes_{E_0}L\), and let a rational domain of
\(\mathcal M(B)\) be defined by \[
             |f_i|\le r_i|g|\quad(1\le i\le s),
             \qquad r_i>0,                              \tag{32}
\] where \(g,f_1,\ldots,f_s\) have no common analytic zero. Empty
systems can instead be handled by the whole affinoid with the harmless
condition \(|1|\le2|1|\). Compactness gives \[
 \delta=\min_{y\in\mathcal M(B)}
             \max\{|g(y)|,r_1^{-1}|f_1(y)|,\ldots,
                                  r_s^{-1}|f_s(y)|\}>0.
                                                               \tag{33}
\] Choose simultaneous finite-stage approximants \(g',f_i'\) in one
\(A_E\) with Banach errors less than \(\delta\min(1,r_1,\ldots,r_s)\).
This is possible by Section 12.1. Their evaluation errors satisfy \[
                 |g'-g|<\delta,
                 |f_i'-f_i|<r_i\delta                 \tag{34}
\] at every point upstairs. Neither vanished nor nilpotent functions are
logarithmed, and none is assumed to have small Banach norm merely
because its spectral norm is zero. They are approximated in the actual
Banach algebra just like the other elements.

Both conjunctions, (32) and the conjunction of its strict versions, are
\textbf{unchanged} by (34). To check all boundary cases, first suppose
\(|g(y)|\ge\delta\). Then \(|g'(y)|=|g(y)|\). For each \(i\) with
\(|f_i(y)|\ge r_i|g(y)|\), (34) forces \(|f_i'(y)|=|f_i(y)|\). For each
\(i\) with \(|f_i(y)|<r_i|g(y)|\), it instead forces
\(|f_i'(y)|<r_i|g'(y)|\). Hence each weak inequality, each equality, and
each strict inequality is preserved at that point. If \(|g(y)|<\delta\),
(33) supplies an \(i\) with \(r_i^{-1}|f_i(y)|\ge\delta\). Its numerator
norm is unchanged, while \(|g'(y)|<\delta\). This inequality remains
strictly false, so both conjunctions remain false. These arguments also
cover \(g(y)=0\) and any number of zero numerators.

The weighted maximum for the approximants is everywhere positive
upstairs. Surjectivity to \(\mathcal M(A_E)\) then shows the
denominator/numerators have no common zero there. They generate the unit
ideal: otherwise a maximal ideal containing them gives an analytic
common zero. Thus they define a genuine rational analytic domain \(R_E\)
by CLD 11.1.5. Its scalar extension has exactly the set (32). Analytic
domain uniqueness, or the universal rational-domain presentation by
adjoining the bounded ratios, identifies it with the given domain
\(R_L\) as an analytic domain. The actual rational localization algebra,
with all nilpotents, is retained. The strict locus
\(O_E=\{|f_i'|<r_i|g'|\ \forall i\}\) is ordinary open in the ambient
affinoid, lies in its rational domain, and has inverse image
\textbf{exactly} the original strict locus \(O_L\). No invariance of
absolute analytic interiors or equality of all topological interiors is
asserted.

The stronger rational-neighbourhood statement of CLD 11.1.5 says that
for any \(y\) and any prescribed ordinary open in an affinoid, one can
choose (32) with \(R_L\) inside that open and \(y\) in its strict locus.
Thus the preceding construction supplies finite-stage ordinary
neighbourhoods of each chosen point, without controlling an entire
infinite scalar fibre. Only finitely many such choices are subsequently
needed on a compact support.

\subsubsection{12.3. Exact finite-stage local
descriptions}\label{exact-finite-stage-local-descriptions}

Let \(\beta\) be ordinary smooth on \(X_L\) near \(y\), and let \(x\) be
its image in \(X\). Choose a downstairs affinoid neighbourhood \(V_0\)
and an ordinary open \(U_0\subset V_0\) containing \(x\). The inverse
image \(U_{0,L}\) is open. Intersect it with an ordinary moment
description open of \(\beta\); on that intersection write
\(\beta=\varphi(t_{f_1},\ldots,t_{f_m})\), with actual analytic units
and a smooth real coefficient defined on an open Euclidean set
containing the coordinate image. This can always be arranged locally by
shrinking the description near \(y\).

Inside \(V_{0,L}\) select a rational neighbourhood \(R_L\) of \(y\)
contained in this intersection, with \(y\) in its strict locus. Section
12.2 gives \(R_L=(R_{E_0})_L\) for a finite stage \(E_0\), and a
descended strict locus. Because \(R_L\subset U_{0,L}\), surjectivity on
\(R_{E_0}\) gives \(R_{E_0}\subset U_{0,E_0}\). Its strict locus is
therefore an ordinary open of the whole \(X_{E_0}\), not just an open of
an exposed affinoid boundary.

On the compact rational affinoid \(R_L\), put \(\delta_j=\min|f_j|>0\).
Apply density in the actual rational affinoid algebra, now viewed over
\(E_0\), to select a common finite \(E\supset E_0\) and
\(u_j\in\mathcal O(R_E)\) with \(\|u_j-f_j\|<\delta_j\) upstairs.
Ultrametricity gives \[
                     |u_j(y')|=|f_j(y')|
                                 \quad(y'\in R_L).     \tag{35}
\] By surjectivity \(u_j\) has no zero on \(R_E\), so it is an actual
unit by the unit criterion. Its inverse exists in that affinoid algebra;
no approximation of an inverse or reduction of the algebra is necessary.
Surjectivity also makes the finite-stage moment image of the \(u_j\)
equal to the upstairs image of the \(f_j\), so the same Euclidean
coefficient \(\varphi\) is defined on it.

On the inverse image \(O_E\) of the strict locus define the ordinary
smooth function \(b_E=\varphi(t_{u_1},\ldots,t_{u_m})\). Then
\(q_E^*b_E=\beta\) on \(O_L=q_E^{-1}O_E\). This is an exact
scalar-section identity. To justify the section claim precisely, CLD
3.4.6 embeds the \((0,0)\) G-form sheaf in numerical functions: on each
compact chart the nullity criterion 3.4.1/3.4.4 says that a zero
numerical coefficient is the zero section; compact charts G-cover the
space. Hence equal numerical scalar sections coincide, and ordinary
scalar sections are included in this sheaf. Alternatively, the combined
moment \((f,u)\) has image in the diagonal \(\{t_f=t_u\}\) by (35); the
difference of the two scalar coefficient functions restricts to zero
there, and moment nullity gives the same conclusion. No universal
moment-dimension equality is used. Smooth differentiation consequently
commutes with this \textbf{exact identity}, giving \(D\beta=q_E^*Db_E\)
on that open.

\subsubsection{12.4. Compact scalar tests and extension by
zero}\label{compact-scalar-tests-and-extension-by-zero}

Now let \(\beta\in A_c^{0,0}(X_L)\) be arbitrary and put
\(F=\operatorname{supp}\beta\). If \(F\) is empty, take \(\alpha=0\) at
any finite stage. Otherwise Section 12.3 gives a neighbourhood
\(O_{i,L}=q_{E_i}^{-1}O_{i,E_i}\) and a descended ordinary function
\(b_{i,E_i}\) near each point of \(F\). Choose finitely many covering
\(F\). Enlarge to a common finite field \(E\subset M\) and base change
all these rational domains, strict loci and unit lists to \(E\).
Associativity (30) shows \(O_{i,L}=q_E^{-1}O_{i,E}\) and
\(\beta=q_E^*b_{i,E}\) there.

The compact set \(F_E=q_E(F)\) is covered by the finitely many
\(O_{i,E}\). Apply Section 4's explicit compact partition, built only
from CLD 15.2.5, to obtain ordinary smooth \(\theta_i\) with compact
support in \(O_{i,E}\) and sum one near \(F_E\). Put \[
                \alpha=\sum_i\widetilde{\theta_i b_{i,E}},
                                                               \tag{36}
\] where tilde means extension by zero to \(X_E\). Each product is
ordinary smooth on its open, with support in the compact
\(\operatorname{supp}\theta_i\subset O_{i,E}\). Since that compact set
is closed in the Hausdorff \(X_E\), every point outside it has an
ordinary zero neighbourhood. Therefore its extension is ordinary smooth
at every exterior point as well. This proves (36) lies in
\(A_c^{0,0}(X_E)\), with no boundary derivative term.

On the inverse image of its defining open each summand pulls back to
\((\theta_i\circ q_E)\beta\). Off that open its pullback is zero, and
\(\theta_i\circ q_E\) is also zero. Thus, \textbf{globally}, \[
 q_E^*\alpha=(\sum_i\theta_i\circ q_E)\beta=\beta.       \tag{37}
\] On \(F\) the sum is one. Outside \(F\) the function \(\beta\) is
zero, so the equality also holds there regardless of that sum. No
descent or saturation of the original support is assumed. No agreement
of the separate \(b_{i,E}\) on downstairs overlaps needs to be presumed;
the finite sum is itself a smooth section, and (37) proves precisely its
required pullback. Its support is contained in a finite union of compact
cutoff supports. The inverse image of that union is compact by Section
12.1. Differentiation commutes with the exact ordinary smooth identity
(37), so every compact upstairs scalar test and its Hessian are covered.
This is stronger than \(C^0\) density and does not use a differentiated
approximation topology at completion.

\subsubsection{12.5. Ambient geometry and local purity after
completion}\label{ambient-geometry-and-local-purity-after-completion}

Goodness follows because inverse images of downstairs affinoid
neighbourhoods are affinoid and contain inverse images of their ordinary
interiors. Strict presentations base change to strict presentations. The
diagonal base change preserves separatedness. Boundarylessness follows
from GR 2.11(5'), as stated in Section 2; ordinary open restrictions
retain it by GR 2.11(3),(4).

We check local purity for \(X_L\), rather than just total dimension.
Around the image of \(y\in X_L\) take a downstairs strict affinoid
neighbourhood \(V_0\) of dimension one. Then \(V_{0,L}\) has dimension
one, giving the upper local bound. Suppose local dimension at \(y\) were
zero. In \(V_{0,L}\) discard its finitely many dimension-one
minimal-prime components, none of which contains \(y\). In the remaining
ordinary open choose a rational affinoid neighbourhood \(R_L\) of \(y\)
inside the inverse image of \(\operatorname{int}_X V_0\), with strict
inequalities at \(y\). The component/domain dimension formula in CLD
11.2.3 gives \(\dim R_L=0\). Section 12.2 descends \(R_L\) to a nonempty
\(R_E\) at a finite field, inside \(X_E\). The actual \(X_E\) is pure
one-dimensional by Sections 3 and 11. Its domain \(R_E\) has dimension
one, including if its chosen real radii make the domain non-strict.
Scalar dimension invariance and analytic domain identification give
\(\dim R_L=\dim(R_E)_L=\dim R_E=1\), a contradiction. Hence every
upstairs point has dimension one. This argument keeps the original
strict ambient space and does not infer an endpoint theorem from a
non-strict radial core.

All top measures below are on the \textbf{actual} pure curves just
proved to exist. Their multiplicities need not remain numerically
unchanged point by point. They are the finite-flat fibre-algebra
dimensions in the canonical calibration. For reference, the usual
generic-length identity is obtained by tensoring a composition series of
the generic Artinian stalk of a chosen downstairs irreducible component
with the flat generic stalk of a chosen upstairs component above it: each
residue-field quotient becomes the generic Artinian algebra of the
base-changed reduction. If its length is \(\mu\) and the original length
is \(\ell\), the length upstairs is \(\mu\ell\). This is CLD 13.2.8 with
\(\mathcal F=\mathcal O_X\); its proof body is in the corrected primary-source supplement.
We use the calibration and the Artinian degree-sum proof of (4), not a
replacement by geometric point counts or a supposed constant extension
degree for an infinite extension.

\subsection{13. Completed forward ascent and actual reverse
descent}\label{completed-forward-ascent-and-actual-reverse-descent}

Assume \(D[h]=0\) downstairs. Given \textbf{any}
\(\beta\in A_c^{0,0}(X_L)\), Section 12.4 supplies finite \(E\) and
\(\alpha\in A_c^{0,0}(X_E)\) with \(\beta=q_E^*\alpha\) exactly. Put
\(h_E=h\circ\pi_E\). Finite algebraic ascent, proved in Sections 3--11,
gives \(D[h_E]=0\). Section 12.5 verifies the pure curve and ambient
hypotheses for both \(X_E\) and \(X_L\). Smooth functoriality and the
\textbf{\(L\)- and \(E\)-calibrated} projection formula (4) yield \[
 \begin{aligned}
 B_{h\circ\pi}(\beta)
 &=\int_{X_L}(h_E\circ q_E)(q_E^*D\alpha)_{\mathrm{Rad},L}\\
 &=\int_{X_E}h_E(D\alpha)_{\mathrm{Rad},E}
 =0.
 \end{aligned}                                             \tag{38}
\] The relevant supports lie in the compact inverse image of
\(\operatorname{supp}\alpha\); the continuous finite multiplier is
bounded there. The signed-measure pushforward formula is therefore
legitimate. Since \(\beta\) was arbitrary, the upstairs zero-current
equation follows. There is no pullback of an arbitrary continuous
current in this proof and no field-degree factor in (38).

Conversely, suppose the upstairs equation holds. For an arbitrary
\(\psi\in A_c^{0,0}(X)\) its pullback has compact support by Section
12.1. Apply smooth functoriality and (4) directly for \(L/K\): \[
             B_h(\psi)=B_{h\circ\pi}(\pi^*\psi)=0.        \tag{39}
\] Thus the downstairs equation holds. This reverse calculation is
included for completeness, not counted as new forward progress.
Ordinary-open localization follows from the ambient checks and the same
arguments on that open; no global quasicompactness or paracompactness is
added. This proves (1) with all its stated scope.

\subsection{14. Comparison, fallibility and
reproduction}\label{comparison-fallibility-and-reproduction}

The controlling P06.2 chapter asks for \textbf{classical pluriharmonic
recognition} of arbitrary continuous zero-current functions on good
separated strict boundaryless spaces, or a counterexample with all those
ambient hypotheses. Neither conclusion is claimed here. The new gate is
unconditional ordinary CLD scalar-current ascent on arbitrary pure
curves under the completion of an arbitrary algebraic extension,
including singular and nonreduced curves and the completed algebraic
closure. It goes beyond merely restating the finite separable curve
draft or imposing scalar stability.

The finite separable dependency has been supplied with its entire
scalar-polarization, root-valuation, tangential-convolution and
calibrated carrier proof. The finite inseparable dependency has been
supplied with its actual tensor, Frobenius and compact scalar-test
proof. Neither predecessor's author audit, publication, native note,
hash, or review verdict is accepted as a mathematical premise. The
finite-stage support and domain gate for completion has been proved
directly. No arbitrary infinite fibre is treated as finite.

In dimension at least two, complementary Hessian tests have positive
bidegree, and the finite separable reconstruction here is only scalar.
Therefore even an extension of the completed-stage argument to all form
degrees would not by itself prove the missing finite differential-form
theorem. Transcendental scalar extensions, universal curve probes,
normalization transport for recognition, and two-sign regularization
remain separate unresolved obligations. Zero opposite currents after
\(L/K\) do not put either continuous sign into the generated CLD-psh
hierarchy. Ordinary smoothness, weak smoothness, real-PL regularity and
GJR harmonicity of \(h\) are not asserted. The catalogue's non-strict
arbitrary-radius invisible boundary example is outside the original
strict boundaryless scope and is not a counterexample supplied here.

The complete frozen mathematical argument, the corrected primary-source
supplement, and the exact-scope review receipt accompany this paper.
The independently audited result is a partial theorem: all original
P06.2 hypotheses and unsolved higher-dimensional and recognition gates
remain as specified above. The native Xean review used
\texttt{gpt-6.1-sol} with \texttt{xhigh} reasoning and returned PASS with
26 external premises and 35 supporting passages. Every premise and
passage was read and its argument/source bindings checked. This audit
is not a formal proof certificate or a priority certification.

The original source appendix selected ranges by newline numbers using
a parser that also counted PDF form feeds. Its inventory was therefore
inaccurate. The historical manuscript and first review are preserved;
a fresh review with correctly matched authentic excerpts governed this
publication. The misleading appendix is replaced by the source audit
below, and its complete corrected quotations accompany the paper.
The finite diagnostic and successful compilation establish checks and
typesetting, respectively, rather than proving the analytic theorem.

\subsection{Bibliography}\label{bibliography}

\begin{itemize}
\tightlist
\item
  A. Chambert-Loir and A. Ducros, \emph{Formes différentielles réelles
  et courants sur les espaces de Berkovich}, arXiv:1204.6277v2, 25 July
  2025. Used statements and proof bodies are specified in Section 2 and
  supplied in the accompanying corrected primary-source supplement. The rational-domain completion lemma is proved in
  Section 12, not attributed to a source slogan.
\item
  W. Gubler and J. Rabinoff, \emph{Classically psh and pluriharmonic
  functions on Berkovich spaces}, arXiv:2506.13548v2, 17 September 2025,
  Section 2.11(3)--(5).
\item
  A. Ducros, \emph{Espaces de Berkovich, polytopes, squelettes et
  th\'eorie des mod\`eles}, Confluentes Math. 4 (2012), no.~4,
  1250007, Theorems 3.2 and 1.2(A),
  \url{https://doi.org/10.1142/S1793744212500077}.
  The compact strict moment-image specialization is used only for
  polyhedrality and the upper dimension bound. The model-theoretic
  foundation is a stated external source dependency.
\item
  Controlling chapter and family dossier: scope data only. Prior finite
  separable and inseparable manuscripts: fallible evidence; their needed
  arguments are reconstructed in this file.
\end{itemize}

\subsection{15. Source audit and corrected supplement}

The authentic corrected packet
\texttt{completed-algebraic-corrected-primary-bodies.md} accompanies
this paper and is embedded directly in the independently reviewed
argument. Its quotations use newline boundaries, with page form feeds
removed solely for display. The main proof in Sections 1--14 is retained
with the recorded printing, component-notation, citation and provenance
corrections. No new mathematical assumption or argument is introduced
by these editorial changes.

The exact external analytic foundations used in the proof are stated
in Section 2 and in the review receipt. They include affinoid
construction, scalar extension and analytic dimension, ordinary form
sheaves, canonical finite-flat calibration and Artinian degrees,
integration and scalar top-measure projection. The completed rational
domain, unit-norm and compact scalar-test descent are proved here.
Ducros 3.2 supplies the compact strict moment-image assertion, with
Theorem 1.2(A) retained as its external model-theoretic foundation.
These foundations are cited and checked; this paper does not claim to
derive all of Berkovich geometry from elementary axioms.

The universal dimension-equality phrase printed in CLD 12.2.4 is unused:
only polyhedrality and the separate upper bound apply. The broad
fixed-index compact-support formulation in CLD 15.2.9 is unused.
Compact cutoffs use CLD 15.2.5 and the direct construction of Section 2.1.
The source's rational-localization algebra, degree and volume-form
subscript printing slips do not alter the constructions supplied in
this proof. Scalar projection has no field-degree factor; restriction
of scalars on the same carrier has the finite factor explicitly
cancelled in the argument. The actual nonreduced algebra is always
retained for integration.

\end{document}
